% 张三慧《大学物理学》第4版第一卷，第8章《狭义相对论基础》
% 主教材PDF物理页283--285（印刷267--269），原习题8.1--8.23。
% 题面忠于原扫描物理条件，少量长段落以结构化中文重述；
% 所有习题是待核验的LaTeX候选，不将思考题8.1--8.11计入。
\section*{第8章\quad 狭义相对论基础}
\subsection*{章末习题8.1--8.23}

\paragraph{习题 8.1}
一根直杆在惯性系 \(S\) 中静止，静止长度为 \(l\)，
与 \(x\) 轴夹角为 \(\theta\)。
系 \(S'\) 沿 \(S\) 的 \(+x\) 轴以速度 \(u\) 匀速运动。
求杆在 \(S'\) 中的长度以及杆与 \(x'\) 轴的夹角。

\paragraph{习题 8.2}
静止时边长为 \(a\) 的正立方体，其一条棱平行运动方向。
该立方体相对于 \(S'\) 系以速度 \(u\) 运动时，
求在 \(S'\) 中测得的体积。

\paragraph{习题 8.3}
\(S\) 系 \(x\) 轴上放一根长尺，两端各安装一支枪；
\(S'\) 系的 \(x'\) 轴上固定另一根有刻度的尺。
两尺从前者旁边经过时，\(S\) 系的观察者\emph{同时}
扣动两枪，在 \(S'\) 尺上击出两个标记。
求 \(S'\) 尺上两标记的刻度间隔；
说明为什么 \(S'\) 观察者对该结果的解释不同。
\textit{题目未给具体尺长及速度，分别记 \(S\) 系两枪距离
\(L\) 和两系相对速度 \(u\)。}

\paragraph{习题 8.4}
在 \(S\) 系中同一地点发生两件事，第二件比第一件
晚 \(2\,\mathrm s\)；而在 \(S'\) 系中，第二件比第一件
晚 \(3\,\mathrm s\)。
求在 \(S'\) 中两事件发生位置之间的距离。

\paragraph{习题 8.5}
在 \(S\) 系中同时发生的两件事相距
\(1\,\mathrm m\)，在 \(S'\) 中两事件的空间距离
为 \(2\,\mathrm m\)。
求 \(S'\) 系中它们的时间间隔。
\textit{默认两事件沿相对运动方向分开。}

\paragraph{习题 8.6}
一艘装备无线电发射与接收器的飞船，
以 \(4c/5\) 的速度远离地球。
宇航员向地球发送信号，经地球反射后，
在飞船钟上过了 \(60\,\mathrm s\) 才收到回波。
求：（1）\emph{地球反射信号时，在飞船系}
测得的地球—飞船距离；
（2）\emph{飞船收到回波时，在地球系}
测得的距离。

\paragraph{习题 8.7}
飞船沿 \(x\) 方向远离地球，
以 \(u=0.80c\) 飞行，宇航员在本船系 \(S'\)
测得超新星爆发事件的坐标为
\[
 t'=-6.0\times10^8\,\mathrm s,\quad
 x'=1.8\times10^{17}\,\mathrm m,\quad
 y'=1.2\times10^{17}\,\mathrm m,\quad z'=0.
\]
假设飞船经过地球时两个系的钟都读零。
求在地球系 \(S\) 中该事件的时空坐标，
以及超新星位置同飞船航行方向的夹角。

\paragraph{习题 8.8}
接上一题：宇航员对超新星坐标的报告
由到达飞船的光线及其方向算出，并非直接测量。
（1）在飞船系 \(S'\) 的什么时刻，
来自超新星的光传到飞船？
（2）宇航员看到超新星后立即把报告发回地球，
地球系的观察者在什么时刻收到该报告？
（3）地球系的观察者直接在什么时刻看到超新星？
\textit{按光信号沿直线以速率 \(c\) 传播，
事件到达飞船/地球都以两系原点世界线判定。}

\paragraph{习题 8.9}
地球系的观察者发现一艘以 \(0.60c\)
向东飞行的宇宙飞船将在 \(5\,\mathrm s\) 后
与一颗以 \(0.80c\) 向西飞行的彗星相撞。
求：（1）飞船中宇航员看到彗星接近的速率；
（2）按宇航员自己的钟，
距离碰撞还剩多长时间。

\paragraph{习题 8.10}
在 \(S'\) 系原点 \(O'\) 发出光线，
光线在 \(S'\) 中与 \(x'\) 轴夹角为 \(\theta'\)。
求光在 \(S\) 中的传播方向，
并证明它在 \(S\) 中的传播速率仍是 \(c\)。
两惯性系沿 \(x/x'\) 方向相对速度 \(u\)。

\paragraph{习题 8.11}
参考图8.9：\(S'\) 沿 \(S\) 的 \(x\) 方向
以恒定速度 \(u\) 运动；一质点在 \(S\)
中沿 \(x\) 运动，速度为 \(v=v(t)\)、加速度为 \(a\)。
证明 \(S'\) 中的加速度
\[
 \boxed{a'=\frac{(1-u^2/c^2)^{3/2}}
 {(1-uv/c^2)^3}\,a}.
\]
讨论当 \(u,v\ll c\) 时这一式子的意义。

\paragraph{习题 8.12}
静质量为 \(m_0\) 的质点，
受沿运动方向的恒力 \(F=F\boldsymbol i\)，
从静止出发。
求经过时间 \(t\) 后的速度 \(v\) 和位移 \(x\)；
再求 \(t\ll m_0c/F\) 与
\(t\gg m_0c/F\) 的两个极限表达式。

\paragraph{习题 8.13}
求：（1）粒子的相对论动量等于非相对论动量两倍
时的速度；（2）相对论动能等于非相对论动能两倍
时的速度。采用同一静质量、同一速度进行比较。

\paragraph{习题 8.14}
北京正负电子对撞机的电子被加速到
动能 \(E_k=2.8\times10^9\,\mathrm{eV}\)。
求：（1）电子速度与光速相差多少
\(\mathrm{m/s}\)；（2）电子的动量；
（3）电子在周长为 \(240\,\mathrm m\)
的储存环内绕行受到的向心力，以及
维持此轨道需多强的垂直磁场。
\textit{采用电子静止能量 \(m_ec^2\simeq0.511\,\mathrm{MeV}\)。}

\paragraph{习题 8.15}
最强的宇宙射线有 \(50\,\mathrm J\) 的能量；
若这束射线由一个质子组成，
求它的速率与光速差多少 \(\mathrm{m/s}\)。

\paragraph{习题 8.16}
给定静质量
\[
 m_p=1.67265\times10^{-27}\,\mathrm{kg},\qquad
 m_n=1.67495\times10^{-27}\,\mathrm{kg},
\]
一个质子与一个中子组成氘核，
静质量
\(m_D=3.34365\times10^{-27}\,\mathrm{kg}\)。
求氘核结合能（MeV）及占氘核静能的百分数。
另给氢原子的电子—质子结合能
\(13.58\,\mathrm{eV}\)，
氢原子静质量
\(m_H=1.67323\times10^{-27}\,\mathrm{kg}\)：
求其结合能占氢原子静能的百分数。

\paragraph{习题 8.17}
太阳主要释放能量的核反应近似为
\[
 4\,{}^1_1\mathrm H\longrightarrow
 {}^4_2\mathrm{He}+2\,{}^0_{+1}\mathrm e.
\]
题给质子质量
\(m_p=1.6726\times10^{-27}\,\mathrm{kg}\)，
氦核质量
\(m_{\rm He}=6.6425\times10^{-27}\,\mathrm{kg}\)，
正电子质量
\(m_e=0.0009\times10^{-27}\,\mathrm{kg}\)。
求：（1）一次反应放出的能量；
（2）质量亏损的效率；
（3）每消耗 \(1\,\mathrm{kg}\) 质子释放能量；
（4）若太阳辐射总功率
\(P=3.9\times10^{26}\,\mathrm W\)，
每秒消耗多少公斤质子；
（5）太阳约含
\(1.5\times10^{30}\,\mathrm{kg}\) 质子，
若持续按此速率消耗可维持多少时间。
\textit{沿用教材的反应式和忽略中微子等附加过程的质量账目。}

\paragraph{习题 8.18}
20世纪60年代发现类星体的极强辐射，
讨论其巨大能量机制的一个可能来源：
物质落入黑洞或中子星释放引力能。
（1）\(1\,\mathrm{kg}\) 物质从远处落到地球表面
释放多少引力能？它占落下物质静能的百分数？
（2）\(1\,\mathrm{kg}\) 物质从远处落到
质量等于太阳、半径为 \(10\,\mathrm{km}\)
的中子星表面，释放多少能量，
其效率与上一题的核反应效率相比如何？
\textit{按教材牛顿引力势能作数量级估算，
不是强引力精确广义相对论解。}

\paragraph{习题 8.19}
两个质子以 \(v=0.5c\) 的速率
从一共同点向相反方向运动。
求：（1）在共同点静止系中每个质子的
动量及能量；（2）在其中一个质子静止系中，
另一个质子的动量与能量。

\paragraph{习题 8.20}
动能为 \(22\,\mathrm{GeV}\) 的电子
轰击静止质子时，
求它们的可资利用能。
\textit{须区分不变量质心总能量
与扣除初态静质量后可用于新粒子形成的能量。}

\paragraph{习题 8.21}
北京正负电子对撞机的电子与正电子
均具有 \(2.8\,\mathrm{GeV}\) 的能量，
相向发生正面碰撞。
求此对撞的可资利用能；
若以同样高速电子轰击静止电子，
为达到相同可资利用能，
所需电子的动能多大？

\paragraph{习题 8.22}
由相对论动量—能量变换证明
\[
 \boxed{E^2-c^2p^2}
\]
为洛伦兹不变量，
即在 \(S,S'\) 两系
\(E^2-c^2p^2=E'^2-c^2p'^2\)。

\paragraph{习题 8.23}
按爱因斯坦质能公式将静质量写为
\(m=E_0/c^2\)，代入牛顿万有引力表达式，
用两个静止质子的静能 \(E_0\) 表示它们之间的引力。
（1）求一个静止质子的静能，单位为 \(\mathrm{GeV}\)；
（2）在现实世界中，两个静止质子间的电力
比引力大约 \(10^{36}\) 倍。
若要使两个\emph{具有相同能量}的质子之间的
引力等于电力，按上述牛顿等效质量模型，
每个质子的能量须达到多少 \(\mathrm{GeV}\)？
这一能量是现今粒子加速器（包括建设中）所述
能量范围 \(10^3\,\mathrm{GeV}\) 的多少倍？
\textit{该模型只是教材要求的数量级推算；
真实极高能引力相互作用须采用适当的场论与时空描述。}
